\documentclass[11pt]{article}
\usepackage[T1]{fontenc}
\usepackage[utf8]{inputenc}
\usepackage{lmodern}
\usepackage[margin=1in]{geometry}
\usepackage{amsmath,amssymb,amsthm,mathtools,booktabs}
\usepackage{microtype}
\usepackage[hidelinks]{hyperref}
\newtheorem{theorem}{Theorem}
\newtheorem{lemma}[theorem]{Lemma}
\newtheorem{corollary}[theorem]{Corollary}
\newcommand{\NN}{\mathbb Z_{\ge0}}
\newcommand{\RR}{\mathbb R_{\ge0}}
\newcommand{\cO}{\mathcal O}
\newcommand{\cI}{\mathcal I}
\DeclareMathOperator{\cone}{cone}
\title{Conditional Schur positivity for binary-word events\\with at most four free positions}
\author{Xiaojian Zeng\\
\small PBC School of Finance, Tsinghua University, 100190\\
\small \texttt{xiaojianzeng@foxmail.com}}
\date{}
\begin{document}
\maketitle

\begin{abstract}
For binary words ending in $1$, mark the ending positions of the consecutive
factors $10$ and $011$. We prove that fixing every event indicator except at
most four positions gives a Schur-positive two-variable antipodal polynomial,
uniformly in word length, the number of fixed events, and the distances between
free positions. Consequently the conditioned counting polynomial is
coefficientwise log-concave in arbitrary weights on the free positions.
The two-position inequality has sharp factor $2$. The antipodal conclusion
also holds for arbitrarily many free positions separated into blocks of size
at most four by adjacent fixed events. The proof reduces arbitrary transfer
matrix words to five separate families and a complete exact certificate of
$695\,500$ parameter coefficients. This is a restricted result for the
binary-word model associated with big descents in three pattern-avoidance
classes; the original log-concavity conjecture remains unresolved by this work.
\end{abstract}

\section{The event model and the conditional theorem}

For $n\ge1$, let $\mathcal W_n$ be the binary words of length $n$ ending in $1$.
For $w=w_1\cdots w_n\in\mathcal W_n$, define
\[
 E(w)=\{i\in\{2,\ldots,n\}:w_{i-1}w_i=10
       \text{ or }(i\ge3\text{ and }w_{i-2}w_{i-1}w_i=011)\}.
\]
The two alternatives cannot occur at the same position. Write
\[
 a_n(A)=\#\{w\in\mathcal W_n:E(w)=A\}
 \qquad(A\subseteq\{2,\ldots,n\}).
\]
Thus $a_n(A)$ specifies the \emph{entire} event set; unattainable sets have
count zero.

This model comes from Elizalde, Rivera, and Zhuang
\cite[Theorems~3.5 and~3.8; Lemmas~3.6 and~3.9]{ERZ}.
A big descent of a permutation $\pi$ is a position $i$ with
$\pi_i>\pi_{i+1}+1$. Their bijections identify the number of big descents in
the avoidance classes
\[
 \{123,132\},\qquad \{123,213\},\qquad \{132,213\}
\]
with $|E(w)|$, and give the corresponding univariate generating function.
Our coordinates are event-ending positions in $w$; we do not identify them
with big-descent positions in $\pi$.
The source conjectures log-concavity of $b_{n,k}(\Pi)$, the number of
$\Pi$-avoiding permutations with $k$ big descents, for all $n\ge0$ and
all $21$ singleton or doubleton subsets $\Pi\subseteq\mathfrak S_3$
\cite[Conjecture~4.6]{ERZ}. The results here do not resolve that assertion,
or even its unrestricted specialization to the three classes above.

For integers $a\ge b\ge0$, the two-variable Schur polynomial is
\begin{equation}\label{eq:schur}
 s_{(a,b)}(x,y)=(xy)^b\sum_{q=0}^{a-b}x^q y^{a-b-q}.
\end{equation}
A symmetric homogeneous polynomial is \emph{Schur-positive} if its expansion
in this basis has nonnegative coefficients; the zero polynomial is included.
For disjoint $S,J\subseteq\{2,\ldots,n\}$, with $m=|J|$, set
\begin{align}
 H_{S,J}(x,y)&=\sum_{T\subseteq J}a_n(S\cup T)
        a_n(S\cup(J\setminus T))x^{|T|}y^{m-|T|},\label{eq:H}\\
 h_j(S,J)&=\sum_{\substack{T\subseteq J\\|T|=j}}
       a_n(S\cup T)a_n(S\cup(J\setminus T)).\label{eq:h}
\end{align}
All event indicators outside $J$ are fixed: positions in $S$ are present
and positions outside $S\cup J$ are absent. Complementation gives
$h_j=h_{m-j}$. With $h_{-1}=0$, equation~\eqref{eq:schur} gives
\begin{equation}\label{eq:delta}
 H_{S,J}=\sum_{j=0}^{\lfloor m/2\rfloor}
       (h_j-h_{j-1})s_{(m-j,j)}.
\end{equation}

\begin{theorem}\label{thm:main}
For every $n\ge1$ and disjoint $S,J\subseteq\{2,\ldots,n\}$ with
$|J|\le4$, the polynomial $H_{S,J}$ is Schur-positive. There is no bound
on $n$, on $|S|$, or on the spacing of $J$.

The same conclusion holds for arbitrary $|J|$ if, in its natural order,
$J$ is partitioned into nonempty blocks of at most four positions and,
between every two successive blocks, there are two adjacent positions
belonging to $S$.
\end{theorem}

\begin{corollary}\label{cor:lc}
Under the first part of Theorem~\ref{thm:main}, define
\begin{equation}\label{eq:G}
 G(t;\boldsymbol z)=\sum_{T\subseteq J}a_n(S\cup T)t^{|T|}
              \prod_{i\in T}z_i=\sum_{k=0}^m g_k(\boldsymbol z)t^k.
\end{equation}
With $g_k=0$ outside $0\le k\le m$, every
$g_k^2-g_{k-1}g_{k+1}$ has nonnegative monomial coefficients.
In particular, the coefficient sequence of $G$ is log-concave for every
nonnegative real assignment to the $z_i$.
\end{corollary}

\begin{corollary}\label{cor:sharp}
For distinct $i,j\notin S$ in $\{2,\ldots,n\}$,
\begin{equation}\label{eq:sharp}
 a_n(S)a_n(S\cup\{i,j\})
 \le2a_n(S\cup\{i\})a_n(S\cup\{j\}).
\end{equation}
The constant $2$ cannot be decreased uniformly.
\end{corollary}

The conditional refinement is the contribution beyond the source's
univariate enumeration. It is related to, but weaker in normalization than,
the antipodal-pairs property of Kahn and Neiman
\cite[equation~(1) and Theorem~1]{KN}. They divide the degree-$j$
antipodal sum by $\binom mj$ and show that their conditional property,
together with absence of internal zeros, implies ultra-log-concavity.
For two free coordinates, writing the four weights as $A,B,C,D$ in
subset order, their hypothesis requires $AD\le BC$.
The actual-word example proving sharpness below has $(A,B,C,D)=(2,1,1,1)$,
so that criterion does not apply. Our result fixes indicators and uses
unnormalized sums; it does not assert conditional negative correlation.

\section{Transfer matrices and arbitrary gaps}

We use the transfer-matrix interpretation of weighted walks
\cite[Section~4.7, Theorem~4.7.1]{StanleyEC1}, with suffix states
$A=$ ending in $0$, $B=$ ending in $01$, and
$C=$ ending in $1$ but not in $01$.
Appending a letter without an event acts on column vectors by $P$;
appending a letter creating an event acts by $Q$, where
\begin{equation}\label{eq:PQ}
 P=\begin{pmatrix}1&0&0\\1&0&0\\0&0&1\end{pmatrix},\qquad
 Q=\begin{pmatrix}0&1&1\\0&0&0\\0&1&0\end{pmatrix}.
\end{equation}
For example, both continuations of state $B$ create an event, whereas
neither continuation of state $A$ does. In state $C$, appending $0$
creates $10$ and appending $1$ is event-free. Thus, with
$v=(1,0,1)^{\mathsf T}$ and $u=(0,1,1)$,
\begin{equation}\label{eq:count}
 a_n(A)=uU_n\cdots U_2v,\qquad
 U_i=\begin{cases}Q&i\in A,\\P&i\notin A.\end{cases}
\end{equation}
The empty product for $n=1$ gives $uv=1$.

The boundary cones needed below are generated by the row vectors
\[
 o_0=(1,0,0),\quad o_1=(0,1,0),\quad
 o_2=(1,0,1),\quad o_3=(0,1,1)
\]
and the column vectors
\[
 r_0=(1,0,0)^{\mathsf T},\quad r_1=(1,1,0)^{\mathsf T},\quad
 r_2=(1,0,1)^{\mathsf T},\quad r_3=(1,1,1)^{\mathsf T}.
\]
Write $\cO=\cone(o_0,o_1,o_2,o_3)$ and
$\cI=\cone(r_0,r_1,r_2,r_3)$, using nonnegative real combinations.

\begin{lemma}\label{lem:cones}
The cones have the descriptions
\begin{align*}
 \cO&=\{(a,b,c):a,b,c\ge0,\ c\le a+b\},\\
 \cI&=\{(a,b,c)^{\mathsf T}:a,b,c\ge0,\ b\le a,\ c\le a\}.
\end{align*}
Right multiplication by $P,Q$ preserves $\cO$, and left multiplication
by $P,Q$ preserves $\cI$.
\end{lemma}
\begin{proof}
For a vector in the displayed row cone, split $c=c_1+c_2$ with
$0\le c_1\le a$ and $0\le c_2\le b$. Its ray coefficients are
$(a-c_1,b-c_2,c_1,c_2)$.
For the column cone choose
$\max(0,b+c-a)\le d\le\min(b,c)$; the coefficients are
$(a-b-c+d,b-d,c-d,d)$. These intervals are nonempty under the respective
inequalities. Conversely every listed ray satisfies those inequalities.
Finally,
\begin{align*}
 (a,b,c)P&=(a+b,0,c),& (a,b,c)Q&=(0,a+c,a),\\
 P(a,b,c)^{\mathsf T}&=(a,a,c)^{\mathsf T},&
 Q(a,b,c)^{\mathsf T}&=(b+c,0,b)^{\mathsf T},
\end{align*}
and these formulas preserve the stated inequalities.
\end{proof}

Write $[m]=\{1,\ldots,m\}$, with $[0]=\varnothing$.
Fix $\ell\in\cO$, $r\in\cI$ and words $M_1,\ldots,M_{m-1}$
in the matrices $P,Q$, allowing empty words. For $T\subseteq[m]$, put
\begin{equation}\label{eq:AT}
 A_T=\ell U_1M_1U_2\cdots M_{m-1}U_mr,
 \qquad U_i=\begin{cases}Q&i\in T,\\P&i\notin T.\end{cases}
\end{equation}
For $m=0$, put $A_\varnothing=\ell r$.
We prove Schur positivity of
\begin{equation}\label{eq:matrixH}
 \mathcal H=\sum_{T\subseteq[m]}A_TA_{[m]\setminus T}x^{|T|}y^{m-|T|}
\end{equation}
when $m\le4$. This implies the first part of Theorem~\ref{thm:main}:
split \eqref{eq:count} at the free positions, listed in reverse order
to match matrix multiplication. The exterior factors give permitted
boundaries by Lemma~\ref{lem:cones}, since $u=o_3$ and $v=r_2$.
Reversing the order preserves cardinalities and complementation.

\begin{lemma}\label{lem:gaps}
Every word in $P,Q$ either has a factorization $ab$ with $a\in\cI$
a column and $b\in\cO$ a row, or belongs to one of the five separate families
\begin{equation}\label{eq:families}
 \cone(I),\quad \cone(P,PQP),\quad \cone(Q,QPQ),\quad
 \cone(PQ,(PQ)^2),\quad \cone(QP,(QP)^2).
\end{equation}
Every word containing $Q^2$ has the former factorization.
\end{lemma}
\begin{proof}
Direct multiplication gives $P^2=P$ and
\begin{equation}\label{eq:rankone}
 Q^2=e_0e_1^{\mathsf T},\qquad
 e_0=(1,0,0)^{\mathsf T}\in\cI,\quad
 e_1^{\mathsf T}=(0,1,0)\in\cO.
\end{equation}
Consequently $LQ^2R=(Le_0)(e_1^{\mathsf T}R)$ has the required factors
by cone invariance, even if one factor vanishes.
Delete repeated $P$'s from any remaining word. If a $QQ$ occurs, use
the same argument. Otherwise the word is alternating, or empty.
The identities
\begin{align}
 PQPQP&=PQP+P,& QPQPQ&=QPQ+Q,\label{eq:odd}\\
 (PQ)^3&=(PQ)^2+PQ,& (QP)^3&=(QP)^2+QP\label{eq:even}
\end{align}
show that each alternating sequence with fixed initial and final letters
satisfies a Fibonacci recurrence as its length increases by two.
Its first two words are exactly the generators in the corresponding
family in \eqref{eq:families}. The empty word is $I$.
All recurrence coefficients are nonnegative, proving the claim for
every length. We do not replace these five families by their joint cone.
\end{proof}

For a gap $M_i=ab$, the scalar in \eqref{eq:AT} factors into two scalars
depending on disjoint free-site sets, with new boundaries $a$ and $b$.
Summing over the two subsets gives
\begin{equation}\label{eq:factorH}
 \mathcal H=\mathcal H_{\mathrm L}\mathcal H_{\mathrm R}.
\end{equation}
Products preserve two-variable Schur positivity: for $a\ge b$ and $c\ge d$,
\begin{equation}\label{eq:product}
 s_{(a,b)}s_{(c,d)}=
 \sum_{i=0}^{\min(a-b,c-d)}s_{(a+c-i,b+d+i)}.
\end{equation}
This follows directly from \eqref{eq:schur}: after removing $(xy)^{b+d}$,
the coefficient of each monomial on either side counts the same possible
splittings between two finite geometric sums. For $m=0$, $\mathcal H=(\ell r)^2$;
for $m=1$, $\mathcal H=(\ell Pr)(\ell Qr)(x+y)$. Both are Schur-positive.
Strong induction therefore handles all rank-one internal gaps once the
separate-family cases are proved.

\section{The finite parameter certificate}

The remaining calculation concerns entire nonnegative parameter cones,
not selected word lengths or sampled parameter values.
Choose one family in \eqref{eq:families} for each gap and write its
generators as $F_{i,0},\ldots,F_{i,d_i-1}$, where $d_i\in\{1,2\}$.
Expand the boundaries and gaps with independent parameters:
\[
 \ell=\sum_{a=0}^3\lambda_a o_a,\qquad
 r=\sum_{b=0}^3\rho_b r_b,\qquad
 M_i=\sum_{c=0}^{d_i-1}\mu_{i,c}F_{i,c}.
\]
Each $A_T$ is multilinear in the $m+1$ groups
$(\lambda,\rho,\mu_1,\ldots,\mu_{m-1})$.
Hence, writing $h_j=\sum_{|T|=j}A_TA_{[m]\setminus T}$,
each $\delta_j=h_j-h_{j-1}$ is homogeneous of degree two in every group.

\begin{lemma}[Exact finite certificate]\label{lem:certificate}
For $m=2,3,4$, every choice of the $m-1$ gap families, and
$0\le j\le\lfloor m/2\rfloor$, every parameter coefficient of
$\delta_j$ is nonnegative.
\end{lemma}
\begin{proof}[Computational proof]
Here is a complete specification of the finite computation.
The group sizes are $\boldsymbol d=(4,4,d_1,\ldots,d_{m-1})$.
For a multilinear index $p=(a,b,c_1,\ldots,c_{m-1})$, calculate the integer
\[
 v_T(p)=o_aU_1F_{1,c_1}U_2\cdots F_{m-1,c_{m-1}}U_mr_b.
\]
A quadratic index $q$ is one unordered pair $(a_g,b_g)$, with
$0\le a_g\le b_g<(\boldsymbol d)_g$, for each parameter group $g$.
Let $\mathcal P(q)$ be the ordered pairs $(p,p')$ of multilinear indices
whose entries in each group form the multiset $\{a_g,b_g\}$.
The coefficient at $q$ in $\delta_j$ is exactly
\begin{equation}\label{eq:coefficient}
 \sum_{(p,p')\in\mathcal P(q)}\left(
 \sum_{|T|=j}v_T(p)v_{[m]\setminus T}(p')
 -\sum_{|T|=j-1}v_T(p)v_{[m]\setminus T}(p')\right),
\end{equation}
where the second sum is empty for $j=0$.
Repeated indices contribute one ordering, and distinct indices contribute
two. Thus no multiplicity or zero coefficient is omitted.

Evaluate \eqref{eq:coefficient} with integer matrix multiplication for
every family choice, every $q$, and every stated $j$.
The exact coefficient arrays supplied in \texttt{coefficients.json}
give Table~\ref{tab:certificate}; all entries are nonnegative.
The table includes the $j=0$ row and zero coefficients.
For each boundary there are $\binom{4+1}{2}=10$ quadratic indices.
Summing over the five choices for one gap gives
$1+4\binom{2+1}{2}=13$ indices. The total number of positions is therefore
\begin{equation}\label{eq:coverage}
 N_m=(\lfloor m/2\rfloor+1)\,100\,13^{m-1}.
\end{equation}
This independently fixes the full domain of the calculation.

\begin{table}[htbp]
\centering
\begin{tabular}{rrrrr}
\toprule
$m$ & Family choices & Coefficient positions & Positive & Negative\\
\midrule
2 & 5   & $2\,600$   & $1\,597$   & 0\\
3 & 25  & $33\,800$  & $15\,633$  & 0\\
4 & 125 & $659\,100$ & $240\,445$ & 0\\
\bottomrule
\end{tabular}
\caption{Complete parameter-coefficient certificate. Every remaining
position is zero.}\label{tab:certificate}
\end{table}

The supplied verifier reconstructs these arrays by a different calculation.
For a homogeneous quadratic polynomial $f$ in one parameter group,
its diagonal coefficient is $f(e_a)$, and its mixed coefficient is
\begin{equation}\label{eq:polarization}
 [\theta_a\theta_b]f=f(e_a+e_b)-f(e_a)-f(e_b)\quad(a<b).
\end{equation}
Hold the other groups formal and apply this identity successively to every
group. The verifier evaluates all Cartesian products of these points,
performs the finite differences, and compares every resulting integer
with the stored array. It also reconstructs the family-choice domain and
checks \eqref{eq:coverage}. It agrees on all $695\,500$ positions.
Both descriptions, \eqref{eq:coefficient} and \eqref{eq:polarization},
recover exact coefficients, rather than infer signs from evaluations.
Nonnegative coefficients imply $\delta_j\ge0$ for every nonnegative
real parameter assignment, proving the lemma.
\end{proof}

\begin{proof}[Proof of Theorem~\ref{thm:main}]
We first prove the matrix statement \eqref{eq:matrixH} by strong induction
for $0\le m\le4$. The cases $m=0,1$ were given above.
For $m\ge2$, a gap with the factorization in Lemma~\ref{lem:gaps}
yields \eqref{eq:factorH}; both factors have
fewer free sites and valid boundary cones, so induction and
\eqref{eq:product} apply. If no such gap occurs, Lemma~\ref{lem:gaps}
places each gap in one of the five families. Lemma~\ref{lem:certificate}
and \eqref{eq:delta} give Schur positivity.
Equation~\eqref{eq:count} then gives the asserted word statement.

For the extension, between each two successive blocks choose adjacent
positions $q,q+1\in S$. These positions lie in an internal gap and
give a $Q^2$ factor there. Cut at all chosen factors using
\eqref{eq:rankone}. Each component has at most four free sites and
boundaries in the required cones. Repeated use of
\eqref{eq:factorH} and \eqref{eq:product} proves the conclusion for
arbitrarily many blocks. No division is involved, including for zero factors.
\end{proof}

\section{Log-concavity, sharpness, and scope}

\begin{proof}[Proof of Corollary~\ref{cor:lc}]
Every monomial that can occur in $g_k^2-g_{k-1}g_{k+1}$ has exponents
$0,1,2$ and total degree $2k$. Let $D$ be its exponent-two positions
and $E$ its exponent-one positions. Set $r=k-|D|$, so $|E|=2r$.
Both factors must contain $D$ and must partition $E$.
For $r\ge1$, the coefficient in question is
\[
 h_r(S\cup D,E)-h_{r-1}(S\cup D,E).
\]
Positions outside $D\cup E$ in the original free set are now fixed absent.
Since $|E|\le |J|\le4$, Theorem~\ref{thm:main} makes this coefficient
nonnegative. For $r=0$ it is the square $a_n(S\cup D)^2$.
Monomials of other degrees or with other exponent patterns have coefficient
zero. This proves every coefficientwise inequality. Evaluation at
nonnegative $\boldsymbol z$ preserves it. Impossible conditionings and
zero weights require no separate normalization.
\end{proof}

\begin{proof}[Proof of Corollary~\ref{cor:sharp}]
For $J=\{i,j\}$, put
$A=a_n(S)$, $B=a_n(S\cup\{i\})$, $C=a_n(S\cup\{j\})$,
and $D=a_n(S\cup\{i,j\})$. Then
\[
 H_{S,J}=AD(x^2+y^2)+2BCxy
        =ADs_{(2,0)}+(2BC-AD)s_{(1,1)},
\]
giving \eqref{eq:sharp}. For $n=4$, $S=\varnothing$ and $J=\{2,4\}$,
the exact word lists are
\[
 \begin{array}{c|cccc}
 T&\varnothing&\{2\}&\{4\}&\{2,4\}\\\hline
 \{w:E(w)=T\}&\{0001,1111\}&\{1001\}&\{0011\}&\{1011\}.
 \end{array}
\]
Thus $(A,B,C,D)=(2,1,1,1)$, which forces the uniform factor to be at least $2$.
\end{proof}

The extension beyond two positions uses additional structure of the word
matrices. To see why the face inequalities alone are insufficient, give the
empty and full subsets of $\{1,2,3\}$ weight $2$, and every other subset
weight $1$. On each of the six two-dimensional faces, $AD=2BC$.
The three-position antipodal coefficients are nevertheless $(4,3,3,4)$,
whose Schur difference is $3-4=-1$. This example concerns general Boolean
weights and is not a counterexample in the binary-word model.

The theorem leaves arbitrary gaps and fixed-event counts unrestricted,
but not arbitrary free-site counts without the stated cuts.
In particular, setting \emph{all} event positions free eventually exceeds
its scope. The four-position cutoff is a proved range, with no claim that
it is optimal. The two-variable polynomial used here is also different
from the quasisymmetric functions of permutation big-descent sets considered
in \cite[Section~4.3]{ERZ}. The comparison with the inspected sources
identifies a conditional refinement; it is not an exhaustive priority claim.

\section{Reproducibility and computational assistance}

The ancillary directory contains \texttt{coefficients.json},
\texttt{verify\_certificate.py}, and a \texttt{README.md}.
Run
\begin{center}
\texttt{python3 manuscript/supplement/verify\_certificate.py}
\end{center}
from the supplied source package, or invoke the script by its absolute path.
It uses Python~3 and NumPy with \texttt{dtype=object}, so all matrix entries,
finite differences, and sign comparisons are arbitrary-precision integers.
No floating-point tolerance or random sample is used.
Besides the coefficient reconstruction, it derives $P,Q$ from suffix
transitions and checks the displayed matrix identities and cone images.

The separation of a stored certificate from reconstruction of its complete
domain also appears in the exact finite-certificate computations of
Zeng, Liu, Ratnavelu, and Huat \cite{ZengCertificates}.
Here the finite objects are parameter monomials, and formula~\eqref{eq:coverage}
ensures complete coverage. No geometric theorem from that work is used.

Direct enumeration additionally checks all $198\,399$ conditionings with
at most four free positions through $n=12$, including $433\,405$ Schur
coefficient positions. It checks $383\,073$ monomial positions in the
log-concavity differences through $n=10$. These finite word checks test
the interpretation and implementation. Uniformity in $n$ follows from
the cone argument, the gap reduction, and the exact polynomial certificate.

OpenAI Codex assisted with proof exploration, literature searches,
code generation, separate proof and implementation checks, and manuscript
preparation. The original coefficient calculation used integer convolution;
a separately written tensor calculation and a later finite-difference
implementation reproduced it. The ancillary verifier is a portable version
of the finite-difference implementation. Separate model contexts do not
constitute external peer review or statistically independent evidence.
The mathematical support is the reduction and reproducible exact arithmetic;
no proof assistant was used.

\bibliographystyle{plain}
\bibliography{references}
\end{document}
